\documentclass[12pt,aspectratio=1610,professionalfont,serif,t]{beamer}
\usepackage[osf,p]{newpxtext}
\usepackage{newpxmath}
\usepackage{fontawesome5}

\usepackage{fancyvrb}
\fvset{%frame=single,
commandchars=\\\{\},
%framesep=1mm,
fontfamily=helvetica,
fontsize=\normalsize
}
%\usepackage[bitstream-charter]{mathdesign}
\usepackage{listings}
\lstset{commentstyle=\color{orange}, keywordstyle=\color{yellow} }


\usetheme{Boadilla}
\useoutertheme{split}
\usecolortheme{albatross}


\setbeamertemplate{blocks}[rounded][shadow=false]
\setbeamersize{text margin left=\dimexpr 1cm+1ex\relax, text margin right=\dimexpr 1cm+1ex\relax}
\addtobeamertemplate{headline}{}{\usebeamerfont{frametitle}\vspace*{1ex}}
\setbeamertemplate{navigation symbols}{}
\setbeamertemplate{page number in head/foot}{}
%\beamerdefaultoverlayspecification{<+->}


\setbeamercolor*{structure}{fg=green!60!black,bg=black}
\setbeamercolor*{normal text}{fg=green!60!black,bg=black}
\setbeamercolor{palette primary}{use={structure,normal text},fg=green,bg=structure.bg}
\setbeamercolor{palette secondary}{use={structure,normal text},fg=structure.fg,bg=structure.bg}
\setbeamercolor{palette tertiary}{use={structure,normal text},fg=structure.fg,bg=structure.bg}
\setbeamercolor{palette quaternary}{use={structure,normal text},fg=white,bg=green!30!black}
\setbeamercolor{palette sidebar quaternary}{use={structure,normal text},fg=white,bg=green!30!black}
\setbeamercolor{title in head/foot}{use={structure,normal text},fg=white,bg=green!30!black}

\setbeamercolor*{example text}{use={structure},fg=green!60!black}
\setbeamercolor*{block body}{use={structure},bg=structure.bg}
\setbeamercolor*{block body alerted}{use={structure},bg=structure.bg}
\setbeamercolor*{block body example}{use={structure},bg=structure.bg}
\setbeamercolor*{block title}{parent=structure,bg=structure.bg}
\setbeamercolor*{block title alerted}{use={structure,alerted text},fg=alerted text.fg,bg=structure.bg}
\setbeamertemplate{navigation symbols}{}
\setbeamertemplate{items}[square]
\setbeamercolor{item projected}{fg=white}
\setbeamercolor*{normal text}{use={structure},fg=white!90!blue,bg=structure.bg}
\setbeamercolor*{separation line}{}
\setbeamercolor*{fine separation line}{}
\setbeamercolor{alerted text}{use={structure},fg=structure.fg}

\usepackage[italian]{babel}
\usepackage[utf8]{inputenc}
\usepackage{pgf}
\usepackage{verbatim}
\usepackage{inconsolata}
\usepackage{listings}
\lstset{language=C, frame=single,
	basicstyle=\ttfamily,
	numbers=left, numberstyle=\tiny\color{gray},
	numbersep=5pt, fancyvrb=true
}

\usepackage[noend]{algorithm2e}
\usepackage{pgf}
\usepackage{tikz}
\usepackage{graphicx}
\usetikzlibrary{%
	arrows,
	arrows.meta,
	positioning,
	calc,
	backgrounds,
	chains,
	matrix,
	patterns,
	automata,
	fit,
	graphs,
	decorations,
	decorations.pathmorphing,
	decorations.pathreplacing,
	decorations.markings,
}

\usepgflibrary{shapes,shapes.geometric}


\usepackage{url}


\usepackage{xmpmulti}
\usepackage{transparent}
\graphicspath{{img/}}

\newcommand{\zero}{\ensuremath{\mathtt{0}}\xspace}
\newcommand{\uno}{\ensuremath{\mathtt{1}}\xspace}
\newcommand{\due}{\ensuremath{\mathtt{2}}\xspace}
\newcommand{\tre}{\ensuremath{\mathtt{3}}\xspace}
\newcommand{\kleene}{\ensuremath{{}^{\ast}}\xspace}
\xspaceaddexceptions{\zero \uno \due \tre \kleene}


\newcommand{\simplebox}[1]{
      \begin{tikzpicture}
              \node[text centered, rectangle, draw=black, thick, text width=\linewidth,align=center, rounded corners, minimum height=1em] (C) {#1};
      \end{tikzpicture}
}
%
% In questo file si mettono i comandi utili solo per questo insegnamento.
%
\author{Gianluca Della Vedova}
\title{Linguaggi e Computabilità}
\institute{Univ. Milano--Bicocca\\
	\texttt{https://www.unimib.it/gianluca-della-vedova}}
\date{\today}



\begin{document}


\begin{frame}[fragile]
	\frametitle{Linguaggi e Computabilità: Lezione 11}
\begin{enumerate}
\item
Proprietà algebriche delle espressioni regolari
\item
Proprietà di chiusura dei linguaggi regolari
\end{enumerate}
\end{frame}

\begin{frame}
\frametitle{Esercizio 1}
Determinare se il linguaggio $L = \{ \zero^{m}  \uno^{n} \; \vert \; m\ge n \ge 1 \}$ è regolare.
In caso affermativo indicare un'espressione regolare che lo denota, altrimenti dimostrare che non è regolare usando il pumping lemma.
% Soluzione: L non è regolare.
% Supponiamo per assurdo che L sia regolare e sia p la costante del pumping lemma.
% Scegliamo z = 0^p 1^p: z appartiene a L (m = n = p >= 1) e |z| = 2p >= p.
% Per ogni decomposizione z = uvw con |uv| <= p e |v| >= 1, uv è un prefisso
% di 0^p, quindi u = 0^j, v = 0^k con k >= 1, w = 0^{p-j-k} 1^p.
% Pompando verso il basso (i = 0): u v^0 w = uw = 0^{p-k} 1^p.
% Poiché k >= 1, il numero di 0 (p-k) è minore del numero di 1 (p),
% quindi m < n e uw non appartiene a L: contraddizione.
% Nota: pompare verso l'alto (i >= 2) non funziona, perché 0^{p+(i-1)k} 1^p
% ha m >= n e appartiene ancora a L.
% Per espandere verso i>=2, prendere il reverse del linguaggio.
\end{frame}

\begin{frame}
\frametitle{Esercizio 2}
Determinare se il linguaggio $L = \{ \zero^{m}  \uno^{n} \; \vert \; m\neq n, m,n \ge 0 \}$ è regolare.
In caso affermativo indicare un'espressione regolare che lo denota, altrimenti dimostrare che non è regolare usando il pumping lemma.
\end{frame}

\begin{frame}
\frametitle{Esercizio 3}
Determinare se il linguaggio $L = \{ \zero^{m}  \uno^{n} \zero^{n-m}\; \vert \; n\ge m \ge 0 \}$ è regolare.
In caso affermativo indicare un'espressione regolare che lo denota, altrimenti dimostrare che non è regolare usando il pumping lemma.
\end{frame}

\end{document}
